\documentclass[10pt,a4paper]{amsart}
\usepackage[margin=19mm]{geometry}
\usepackage{amsmath,amssymb,mathtools}
\usepackage[scr=rsfs,cal=euler]{mathalfa}
\usepackage{xfrac}
\usepackage[unicode,hidelinks,bookmarks=false]{hyperref}
\newcommand{\ie}{i.e.\ }
\newcommand{\cf}{cf.\ }
\newcommand{\resp}{respectively\ }
\newcommand{\Subsection}{\subsection}
\providecommand{\nobreakdash}{\nobreak\hbox{-}}
\setlength{\emergencystretch}{2em}
\multlinegap0pt
\usepackage{mathtools}
\usepackage{amssymb}
\usepackage{xcolor}
\usepackage{enumerate}
\usepackage[shortlabels]{enumitem}
\usepackage{esint}
\newtheorem{lemma}{Lemma}
\DeclareMathOperator{\dist}{dist}
\newcommand{\eps}{\varepsilon}
\DeclareMathOperator{\supp}{supp}
\newcommand{\xeps}{\left(\frac{x}{\varepsilon}\right)}
\title{Quantitative periodic homogenization of\\ parabolic equations with large drift and potential}
\author{Kshitij Sinha (IIT Bombay)}
\date{Source: arXiv:2509.22003v3 (CC BY 4.0)}
\begin{document}
\maketitle

\noindent\emph{Source and licence.} Quantitative periodic homogenization of\\ parabolic equations with large drift and potential. Version arXiv:2509.22003v3 (CC BY 4.0), distributed by arXiv under the Creative Commons Attribution 4.0 International licence, http://creativecommons.org/licenses/by/4.0/, as recorded in the arXiv licence metadata for this exact version and shown on the abstract page https://arxiv.org/abs/2509.22003v3. Full licence notice: \texttt{licenses/arxiv-2509.22003-CC-BY-4.0.txt}.\par\smallskip

\noindent\textit{Editorial scope.} Counted unit: the article's lemma lem:bound-weak-form-2 (source0.tex lines 1496--1506). The article's complete original proof of it is delivered with it (source0.tex lines 1507--1580, one \texttt{proof} environment of 74 lines). No mathematics is written, repaired or re-derived: the statement and the proof are the article's own lines, in the article's own notation, and no retained line is substituted or re-flowed.  Retained uncounted as the source-local context the proof needs: lines 889--891; lines 999--1010; lines 1376--1380; lines 1477--1493; lines 1614--1621; lines 1633--1640; lines 1642--1651.  Nothing was omitted from any retained span, so the package carries no omission record.  No retained line contains a citation, so the wrapper carries no bibliography and no bibliography entry.  No retained line contains a figure environment, an image-inclusion command or drawing code, so no image is delivered and the package declares no asset dependency.  Editorial formatting only, in addition: an AMS \texttt{amsart} preamble that restates the article's own declarations and macros used by the retained lines, namely the environments \texttt{lemma} on separate counters, together with the standard packages those lines need.  The retained environments print in this document as Lemma 1 in order of appearance, whereas the article prints the same items with the numbering of its own full document.  The complete native source of the article as distributed in the arXiv e-print for this exact version is bundled unchanged as \texttt{sources/185781/source0.tex}.
\begin{equation}\label{eq:defn-weps}
w_\eps := f_\eps - f_0 - \eps \omega\left(\frac{x}{\eps}\right)\cdot S_\eps \left( \eta_1^\eps \eta_2^\eps \nabla f_0 \right),
\end{equation}

\begin{lemma}\label{lem:wkform-bound}
Let $w_\eps$ be defined by \eqref{eq:defn-weps}. Then for any $\psi\in\mathrm L^2(0,T;\mathrm H^1_0(\Omega))$, we have
\begin{equation}\label{eq:wkform-bound}
\begin{aligned}
\int_0^T & \left\langle \zeta_\eps\partial_{t}w_{\varepsilon},\psi \right\rangle_{\mathrm{H}^{-1}(\Omega) \times \mathrm{H}_0^1(\Omega)} 
+ \int_{\Omega_T} \Theta_\eps\nabla w_{\varepsilon} \cdot \nabla \psi 
\\
& \le C \eps^{\frac14} \left\{ \left\Vert \partial_tf_0\right\Vert_{L^2(0,T;W^{1,d}(\Omega))} + \left\Vert f_0\right\Vert_{\mathrm L^2(0,T;\mathrm H^2(\Omega))} + \left\Vert f_{{\rm{in}}}\right\Vert_{L^2(\Omega)} \right\}
\left\Vert \psi\right\Vert_{\mathrm L^2(0,T;\mathrm H^1_0(\Omega))}
\end{aligned}
\end{equation}
 \end{lemma}

\begin{equation}\label{eq:defn-zeps}
z_\eps := f_\eps - f_0 
- \eps \omega\left(\frac{x}{\eps}\right)\cdot S_\eps \left( \eta_3^\eps \eta_4^\eps \nabla f_0 \right)
- \eps^2 \sum_{i,j=1}^d \frac{\mathfrak{B}_{(d+1)ij}\xeps}{\zeta\xeps} \frac{\partial}{\partial x_i} \left(S_\eps \left( \eta_3^\eps \eta_4^\eps \frac{\partial f_0}{\partial x_j}\right)\right),
\end{equation}

\begin{equation}\label{eq:wk-form-2}
\begin{aligned}
\int_0^T & \left\langle \zeta_\eps\partial_{t}z_{\varepsilon},\psi \right\rangle_{\mathrm{H}^{-1}(\Omega) \times \mathrm{H}_0^1(\Omega)} 
+ \int_{\Omega_T} \Theta_\eps\nabla z_{\varepsilon} \cdot \nabla \psi
= - \int_{\Omega_T} (\zeta_\eps -1)\partial_{t}f_{0}\psi
\\
& 
+ \int_{\Omega_T}\left(\Theta_h - \Theta_\eps\right) \left( \nabla f_0 - S_\eps \left( \eta_3^\eps \eta_4^\eps \nabla f_0 \right) \right) \cdot \nabla \psi
- \eps \int_{\Omega_T} \Theta_\eps \nabla \left\{S_\eps \left( \eta_3^\eps \eta_4^\eps \nabla f_0 \right)\right\} \omega_\eps \cdot \nabla \psi
\\
&- \eps^2 \int_{\Omega_T} \Theta_\eps\nabla \left(\sum_{i,j=1}^d \frac{\mathfrak{B}_{(d+1)ij}\xeps}{\zeta\xeps} \frac{\partial}{\partial x_i} \left(S_\eps \left( \eta_3^\eps \eta_4^\eps \frac{\partial f_0}{\partial x_j}\right)\right)\right) \cdot \nabla \psi
\\
& - \eps^2 \sum_{j,k=1}^d \int_{\Omega_T} \mathfrak{B}_{k(d+1)j}\xeps \partial_t \left\{S_\eps \left( \eta_3^\eps \eta_4^\eps \frac{\partial f_0}{\partial x_j} \right)\right\} \frac{\partial\psi}{\partial x_k}
\\
& - \eps \sum_{i,j,k=1}^d \int_{\Omega_T} \mathfrak{B}_{kij}\xeps \frac{\partial}{\partial x_i} \left( S_\eps \left( \eta_3^\eps \eta_4^\eps \frac{\partial f_0}{\partial x_j} \right)\right) \frac{\partial\psi}{\partial x_k}.
\end{aligned}
\end{equation}

\begin{lemma}\label{smooth: compact support}
For any function $\phi \in \mathrm{L}^{2}(\Omega_T)$ with $\supp{\phi} \subset \subset \Omega_T$, we have the following
\begin{enumerate}
\item $|| S_{\eps}(\phi)||_{\mathrm{L}^2(\Omega_T)} \leq C ||\phi||_{\mathrm{L}^2(\Omega_T)}$
\item $\eps || \nabla S_{\eps}(\phi)||_{\mathrm{L}^2(\Omega_T)} + \eps^2 || \nabla^2 S_{\eps}(\phi)||_{\mathrm{L}^2(\Omega_T)} \leq C ||\phi||_{\mathrm{L}^2(\Omega_T)}$
\end{enumerate}
where the constant $C$ depends upon dimension $d$.
\end{lemma}

\begin{lemma}\label{smooth: extension}
For $Y$-periodic function $g$ and $f \in \mathrm{L}^{p}(\Omega_T)$ consider any subset $\omega \subset \subset \Omega$ and $0<s<\tau<T$ then we have the following
\begin{enumerate}
\item $\left\Vert g\xeps S_{\eps}(f) \right\Vert_{\mathrm{L}^p((s,\tau)\times \omega)} \leq C ||g||_{\mathrm{L}^p(Y)} || f ||_{\mathrm{L}^p((s-\eps^2, \tau + \eps^2)\times\omega(\eps))}$
\item $\eps\left\Vert g\xeps \nabla S_{\eps}(f) \right\Vert_{\mathrm{L}^p((s,\tau)\times \omega)} \leq C ||g||_{\mathrm{L}^p(Y)} || f ||_{\mathrm{L}^p((s-\eps^2, \tau + \eps^2)\times\omega(\eps))}$
\end{enumerate}
where $\omega(\eps):=\{x\in \Omega: \dist(x,\omega)<\eps\}$, $0<\eps <\min\{s,\tau\}$ and the constant $C$ depends upon dimension $d$.
\end{lemma}

\begin{lemma}\label{lemma: quotient}
For $g \in \mathrm H^1(\Omega) \cap \mathrm L^{\infty}(\Omega)$ satisfying the following bound
\[
0 < a \leq g(x) \leq b \qquad \text{ for a.e. } x \in \Omega
\]
we get $\frac{1}{g} \in \mathrm H^1(\Omega) \cap \mathrm L^{\infty}(\Omega) $ and 
\[
\nabla \left[\frac{1}{g}\right] = - \frac{1}{g^2} \nabla g
\]
\end{lemma}

\begin{lemma}\label{lem:bound-weak-form-2}
Let $z_\eps$ be defined by \eqref{eq:defn-zeps}. Then for any $\psi \in \mathrm{L}^2(0,T;\mathrm H_0^1(\Omega))$, we have
\begin{equation*}
\begin{aligned}
\int_0^T & \left\langle \zeta_\eps\partial_{t}z_{\varepsilon},\psi \right\rangle_{\mathrm{H}^{-1}(\Omega) \times \mathrm{H}_0^1(\Omega)} 
+ \int_{\Omega_T} \Theta_\eps\nabla z_{\varepsilon} \cdot \nabla \psi
\\
& \le C \eps^{\frac12} \left\{ \left\Vert \partial_tf_0\right\Vert_{L^2(0,T;W^{1,d}(\Omega))} + \left\Vert f_0\right\Vert_{\mathrm L^2(0,T;\mathrm H^2(\Omega))} + \left\Vert f_{{\rm{in}}}\right\Vert_{L^2(\Omega)} \right\} \left\Vert \psi \right\Vert_{\mathrm L^2(0,T;\mathrm H^1_0(\Omega))}.
\end{aligned}
\end{equation*}
\end{lemma} 

\begin{proof}
Thanks to the expression in \eqref{eq:wk-form-2}, we obtain 
\begin{equation*}
\begin{aligned}
\int_0^T & \left\langle \zeta_\eps\partial_{t}z_{\varepsilon},\psi \right\rangle_{\mathrm{H}^{-1}(\Omega) \times \mathrm{H}_0^1(\Omega)} 
+ \int_{\Omega_T} \Theta_\eps\nabla z_{\varepsilon} \cdot \nabla \psi
\le \left\vert \int_{\Omega_T} (\zeta_\eps -1)\partial_{t}f_{0}\psi \right\vert
\\
& 
+ \left\vert \int_{\Omega_T}\left(\Theta_h - \Theta_\eps\right) \left( \nabla f_0 - S_\eps \left( \eta_3^\eps \eta_4^\eps \nabla f_0 \right) \right) \cdot \nabla \psi \right\vert
+ \eps \left\vert \int_{\Omega_T} \Theta_\eps \nabla \left\{S_\eps \left( \eta_3^\eps \eta_4^\eps \nabla f_0 \right)\right\} \omega_\eps \cdot \nabla \psi \right\vert
\\
& + \eps^2 \left\vert \int_{\Omega_T} \Theta_\eps\nabla \left(\sum_{i,j=1}^d \frac{\mathfrak{B}_{(d+1)ij}\xeps}{\zeta\xeps} \frac{\partial}{\partial x_i} \left(S_\eps \left( \eta_3^\eps \eta_4^\eps \frac{\partial f_0}{\partial x_j}\right)\right)\right) \cdot \nabla \psi \right\vert
\\
& + \eps^2 \left\vert \sum_{j,k=1}^d \int_{\Omega_T} \mathfrak{B}_{k(d+1)j}\xeps \partial_t \left\{S_\eps \left( \eta_3^\eps \eta_4^\eps \frac{\partial f_0}{\partial x_j} \right)\right\} \frac{\partial\psi}{\partial x_k} \right\vert
\\
& + \eps \left\vert \sum_{i,j,k=1}^d \int_{\Omega_T} \mathfrak{B}_{kij}\xeps \frac{\partial}{\partial x_i} \left( S_\eps \left( \eta_3^\eps \eta_4^\eps \frac{\partial f_0}{\partial x_j} \right)\right) \frac{\partial\psi}{\partial x_k} \right\vert =: \ell_1 + \ell_2 + \ell_3 + \ell_4 + \ell_5 + \ell_6.
\end{aligned}
\end{equation*}
Note that $\ell_1$ is the same as $I_1$ in the proof of Lemma \ref{lem:wkform-bound}. Hence we have
\begin{equation}\label{eq:bd-ell-1}
\ell_1 \le C \eps\left\Vert \partial_t f_0 \right\Vert_{L^2(0,T;\mathrm W^{1,d}(\Omega))} \left\Vert \psi\right\Vert_{\mathrm L^2(0,T;\mathrm H^1_0(\Omega))} .
\end{equation}
Also note that $\ell_2$ is similar to $I_3$ in the proof of Lemma \ref{lem:wkform-bound}, except for the choice of the cut-off functions. Mimicking the arguments that dealt with $I_3$ yields the following bound:
\begin{equation}\label{eq:bd-ell-2}
\ell_2 \le C \eps^{\frac12}\left\{
\left\Vert f_0 \right\Vert_{\mathrm L^2(0,T;H^2(\Omega))}
+ \left\Vert \partial_t f_0 \right\Vert_{\mathrm L^2(0,T;W^{1,d}(\Omega))}
+ \left\Vert  f_{\text{in}} \right\Vert_{\mathrm L^2(\Omega)}
\right\} \left\Vert \nabla \psi \right\Vert_{\mathrm L^2(\Omega_T)}.
\end{equation}
Further, $\ell_3$ is similar to $I_4$ in the proof of Lemma \ref{lem:wkform-bound} (again, except for the choice of the cut-off functions). A repeat of those arguments results in
\begin{equation}\label{eq:bd-ell-3}
\ell_3 \le C \eps^{\frac12}
\left\Vert f_0\right\Vert_{\mathrm L^2(0,T;\mathrm H^2(\Omega))}
\left\Vert \nabla \psi \right\Vert_{\mathrm L^2(\Omega_T)}.
\end{equation}
Note that
\begin{equation}\label{eq:inter-1}
\frac{\partial}{\partial x_i} \left(S_\eps \left( \eta_3^\eps \eta_4^\eps \frac{\partial f_0}{\partial x_j}\right)\right)
= S_\eps \left( \frac{\partial}{\partial x_i}\left(\eta_3^\eps \eta_4^\eps\right) \frac{\partial f_0}{\partial x_j}\right)
+ S_\eps \left( \eta_3^\eps \eta_4^\eps \frac{\partial^2 f_0}{\partial x_i \partial x_j}\right)
\end{equation}
and
\begin{equation}\label{eq:inter-2}
\nabla^2 S_\eps\left(\eta_3^\eps \eta_4^\eps \nabla f_0\right) 
= \nabla S_\eps\left(\nabla\left(\eta^\eps_3 \eta^\eps_4\right) \nabla f_0\right) 
+ \nabla S_\eps\left(\eta^\eps_3 \eta^\eps_4 \nabla^2 f_0\right)
\end{equation}
Considering $\ell_4$, we observe that
\begin{equation*}
\begin{aligned}
\ell_4 & = \eps \left\vert \int_{\Omega_T} \Theta_\eps \left(\sum_{i,j=1}^d \nabla \left(\frac{\mathfrak{B}_{(d+1)ij}}{\zeta}\right)\xeps \frac{\partial}{\partial x_i} \left(S_\eps \left( \eta_3^\eps \eta_4^\eps \frac{\partial f_0}{\partial x_j}\right)\right)\right) \cdot \nabla \psi \right\vert
\\
& \qquad + \eps^2 \left\vert \int_{\Omega_T} \Theta_\eps \left(\sum_{i,j=1}^d \frac{\mathfrak{B}_{(d+1)ij}\xeps}{\zeta\xeps} \nabla \left( \frac{\partial}{\partial x_i} \left(S_\eps \left( \eta_3^\eps \eta_4^\eps \frac{\partial f_0}{\partial x_j}\right)\right)\right)\right) \cdot \nabla \psi \right\vert.
\end{aligned}
\end{equation*}
Using \eqref{eq:inter-1} and \eqref{eq:inter-2} and a few properties of the smoothing operator $S_\eps$ (see Lemma \ref{smooth: compact support} and \ref{smooth: extension} in the Appendix), we obtain the bound:
\begin{equation}\label{eq:bd-ell-4}
\begin{aligned}
\ell_4 & \leq C  \Vert \nabla f_0 \Vert_{\mathrm L^2(\Omega_{T,3\eps})} \Vert \nabla \psi \Vert_{\mathrm L^2(\Omega_{T,3\eps})} 
+ C \eps \Vert \nabla^2 f_0 \Vert_{\mathrm L^2(\Omega_T)} \Vert \nabla \psi \Vert_{\mathrm L^2(\Omega_T)}
\\
& \le C \eps^\frac12 \left\Vert \nabla^2 f_0 \right\Vert_{\mathrm L^2(\Omega_T)} \Vert \nabla \psi \Vert_{\mathrm L^2(\Omega_T)}.
\end{aligned}
\end{equation}
Note that to arrive at the estimate \eqref{eq:bd-ell-4}, we have used the fact that $\frac{\mathfrak{B}_{(d+1)ij}}{\zeta}\in \mathrm H^1(\Omega)$ -- see Lemma \ref{lemma: quotient} in the Appendix. We emphasise here that it is at this point we require the coefficient $\zeta$ to be more regular. The terms $\ell_5$ and $\ell_6$ are similar to $I_5$ in the proof of Lemma \ref{lem:wkform-bound} (again, except for the choice of the cut-off functions) and hence can be handled in a similar fashion while taking into account the identity:
\begin{equation*}
\partial_t \left( S_\eps\left(\eta^\eps_3 \eta^\eps_4 \nabla f_0\right) \right) 
= S_{\eps}(\partial_t(\eta^\eps_3 \eta^\eps_4) \nabla f_0) 
+ \nabla S_{\eps}\left(\eta^\eps_3 \eta^\eps_4\partial_t f_0\right) - S_\eps\left(\nabla \left(\eta^\eps_3 \eta^\eps_4\right) \partial_t f_0\right)
\end{equation*}
As the arguments leading to the rest of the  desire estimates are a repeat of those made earlier, we are skipping the details.
\end{proof}

\end{document}
